\documentclass[screen]{beamer}
\usepackage[T1]{fontenc}
\usepackage[utf8]{inputenc}
\usepackage{color}
\usepackage{graphicx}
\graphicspath{{images/presentation/}}
\usepackage[absolute,overlay]{textpos}
\usepackage{listings}
\lstset{
  language=Python,
  backgroundcolor=\color{gray},
  keywordstyle=\bfseries\ttfamily\color[rgb]{0,0,1},
  identifierstyle=\ttfamily,
  commentstyle=\color[rgb]{0.133,0.545,0.133},
  stringstyle=\ttfamily\color[rgb]{0.627,0.126,0.941},
  showstringspaces=false,
  basicstyle=\small,
  stepnumber=1,
  numbersep=10pt,
  tabsize=2,
  breaklines=true,
  prebreak = \raisebox{0ex}[0ex][0ex]{\ensuremath{\hookleftarrow}},
  breakatwhitespace=false,
  aboveskip={1.5\baselineskip},
  columns=fixed,
  upquote=true,
  extendedchars=true,
  inputencoding=utf8
}
%
%   For at dette skal funke:
%   wget http://www.ntnu.no/selvhjelpspakken/ntnubeamer/beamerntnu.tar.gz
%   extract til: /usr/share/texmf/tex/latex/beamer/themes
%   kjør sudo mktexlsr
%   sourcE: http://www.ntnu.no/selvhjelpspakken/ntnubeamer/

% Bruk NTNU-temaet for beamer (her i bokmålvariant), alternativer er
% ntnunynorsk og ntnuenglish.
\usetheme{ntnuenglish}

% Angi tittelen, vi gir også en kortere variant som brukes nederst på
% hver slide:
\title[SHA-3 in Tahoe-LAFS]%
{Experimenting with SHA-3 candidates in Tahoe Least Authorithy File System}

% Denne kan du også bruke hvis det passer seg:
%\subtitle{Valgfri undertittel}

% Angir foredragsholder, også en (valgfri) kortversjon i
% hakeparanteser først som kommer nederst på hver slide:
\author[EH \& PR]{Eirik Haver and Pål Ruud}

% Institusjon. Bruk gjerne disse slik det passer best med det du vil
% ha.  Valgfri kortversjon her også
\institute[NTNU]{Department of Telematics}

% Datoen blir også trykket på forsida.
\date{December 7th 2010}
%\date{} % Bruk denne hvis du ikke vil ha noe dato på forsida.


% Fra her av begynner selve dokumentet
\begin{document}

% Siden NTNU-malen har en annen bakgrunn på forsida, må dette gjøres
% i en egen kommando, ikke på vanlig beamer-måte:
\ntnutitlepage

% Her begynner første slide/frame, (nummer to etter forsida).
\begin{frame}
  \frametitle{Goal}

  Measure and compare performance of the current SHA-3 candidates within the
  Tahoe-Least Authority File System.

\end{frame}

\begin{frame}
  \begin{textblock*}{76mm}[1,1](83mm,18mm)
  \includegraphics[scale=0.4]{tahoe-logo.pdf}
  \end{textblock*}

  \begin{itemize}
  \item Distributed FS
  \item Principle of Least Authority
  \item Erasure coding, k-out-of-n
  \item Written in Python
  \item Uses (double) SHA256 today
    \begin{itemize}
        \item Because of Length-Extension-Attack
    \end{itemize}
  \end{itemize}
\end{frame}

\begin{frame}
  \frametitle{NIST SHA-3 Competition}
  \begin{itemize}
  \item Select SHA-3 implementation by 2nd Quarter 2012
  \item Currently 14 candidates, one from NTNU
  \item Resistant to Length-Extension-Attacks
  \end{itemize}
\end{frame}

\begin{frame}
  \frametitle{What we have done}
  \begin{enumerate}
  \item Made Python-binding for all candidates (C implementations)
  \item Modified Tahoe-LAFS to use the bindings
  \item Benchmarked Tahoe-LAFS with the different candidate implementations
  \end{enumerate}
\end{frame}

\begin{frame}
  \frametitle{Python-bindings}
  \begin{itemize}
  \item Optimized version of candidates available at eBACS: ECRYPT Benchmarking
  of Cryptographic Systems. (SUPERCOP toolkit)
  \item We have used the fastest, capable C-implementations available for our target
  architecture, 32-bit.
  \item Made them available in Python by the use of Cython in ``SHA3Lib''
  \end{itemize}
\end{frame}

\begin{frame}[fragile]
  \frametitle{Modification of Tahoe-LAFS}
  \begin{enumerate}
  \item Quick hack to quickly swap used hash implementation
  \item Benchmarking data (number of hash operations, timers)
  \end{enumerate}

  \begin{lstlisting}
    from sha3lib import bmw256 as SHA256
  \end{lstlisting}

\end{frame}

\begin{frame}
    \frametitle{Tahoe-LAFS is SLOW}
    \begin{itemize}
        \item On 1 Gbit network - with modern computers
        \item Tahoe-LAFS output is about 10-12 Megabyte/s
        \item SCP output is about 80 Megabyte/s
        \item High-CPU load on client but not on storage nodes
    \end{itemize}
\end{frame}

\begin{frame}
    \frametitle{Benchmarking}

    INSERT graphs.. % TODO
\end{frame}

\begin{frame}
    \frametitle{Criticism}
    \begin{itemize}
        \item SUPERCOP is a good idea, however we see problems
            \begin{enumerate}
                \item Speed of Update() is neglected
                \item Implementations are not correct according to specification
                \item Some implementations fails to fulfill the basic requirements
                of HASH functions
            \end{enumerate}
        \item Our implementations are not the fastest possible
            \begin{enumerate}
                \item Some performance loss by using Cython (?)
                \item Not (always) the fastest versions from SUPERCOP
            \end{enumerate}
    \end{itemize}
\end{frame}

\begin{frame}
    \frametitle{Remaining}
    \begin{enumerate}
        \item Figure out inconsistencies.
        \item Finish report....
    \end{enumerate}
\end{frame}

\begin{frame}
    \frametitle{Preliminary Conclusion}
    \begin{enumerate}
        \item There is a performance gain by switching to another hash
        function.
        \item However, the gain is small and we believe improving other areas
        could improve Tahoe performance more.
    \end{enumerate}
\end{frame}

\end{document}
